% ======================================================================
\section{Background and the setup}\label{sec:setup}
% ======================================================================

In this section we fix the conventions, introduce the light-cone Hamiltonian of pure Yang--Mills theory and its eikonal form in the presence of fast (valence) modes, and review the light-cone perturbation theory that is used to construct the soft gluon wave function of the projectile. At the end of the section the unitary operator $\Omega$, which generates the soft gluon wave function from the valence state, is defined and parametrised by a set of coefficient functions. These coefficients are computed in Section~\ref{sec:omega}.

\subsection{Light-cone coordinates and conventions}\label{sec:conventions}

We review the basic setup of QCD on the light cone~\cite{Lepage:1980fj,Brodsky:1997de}. In light-cone coordinates a four-vector $x^\mu$ has the components
\begin{align}
x^\mu=(x^+,x^-,\v x)\,,\qquad x^\pm=\frac{1}{\sqrt2}\,(x^0\pm x^3)\,,\qquad \v x=(x^1,x^2)\,.
\end{align}
Quantisation is carried out on the surfaces of fixed $x^+$. Hence $x^+$ plays the role of the light-cone time, along which states evolve, and $x^-$ is the longitudinal coordinate. Momenta are written in the same way,
\begin{align}
p^\mu=(p^+,p^-,\v p)\,,\qquad p^\pm=\frac{1}{\sqrt2}\,(p^0\pm p^3)\,,\qquad p\cdot x=p^+x^-+p^-x^+-\v p\cdot\v x\,.
\end{align}
The projectile moves in the $+$ direction. The longitudinal momentum $p^+$ is positive for a physical particle, and $p^-$ is the light-cone energy, the generator of translations in $x^+$. For an on-shell massless particle $p^2=2p^+p^--\v p^2=0$, so that
\begin{align}\label{eq:energy}
p^-=\frac{\v p^2}{2p^+}\,.
\end{align}
All energy denominators of light-cone perturbation theory are differences of such light-cone energies. We use the abbreviations
\begin{align}\label{eq:measure}
&d^3p\equiv dp^+\,d^2p\,,\qquad \delta^{(3)}(k-p)\equiv\delta(k^+-p^+)\,\delta^{(2)}(\v k-\v p)\,,\qquad \int_x\equiv\int d^2x\,,\notag\\
&E_g(p)\equiv\frac{\v p^2}{2p^+}\,,\qquad E_{gg}(p,q)\equiv\frac{\v p^2}{2p^+}+\frac{\v q^2}{2q^+}\,,
\end{align}
where a momentum argument without boldface, as in $E_g(p)$ or $a(p)$, stands for the pair $(p^+,\v p)$. Where no confusion is possible, transverse vectors are also written without boldface, and $\v p^2$ always denotes the square of the transverse vector. Transverse indices are $i,j,k,l,m=1,2$; since the transverse metric is Euclidean, they are written as upper or lower indices without distinction, and repeated indices are summed. Colour indices in the adjoint representation are $a,b,c,d,e$, and $f^{abc}$ are the structure constants of $SU(N_c)$. The natural logarithm is denoted by $\log$.

We work in the light-cone gauge $A^+=0$. In this gauge the transverse components $A^a_i$, $i=1,2$, are the only dynamical degrees of freedom; $A^-$ is not dynamical and is expressed through the transverse components by the equations of motion. The gluon creation and annihilation operators satisfy
\begin{align}\label{eq:aadagger}
\big[a^a_i(k),a^{\dagger b}_j(p)\big]=(2\pi)^3\,\delta^{ab}\,\delta_{ij}\,\delta^{(3)}(k-p)\,,
\end{align}
with all other commutators vanishing. For the Fourier transforms in the transverse plane we use
\begin{align}\label{eq:FTconventions}
&a(k^+,\v k)=\int d^2z\,e^{-i\v k\cdot\v z}\,a(k^+,\v z)\,,\qquad
a^\dagger(k^+,\v k)=\int d^2z\,e^{i\v k\cdot\v z}\,a^\dagger(k^+,\v z)\,,\notag\\
&f(\v k)=\int d^2x\,e^{i\v k\cdot\v x}\,f(\v x)\,,\qquad f(\v x)=\int\frac{d^2k}{(2\pi)^2}\,e^{-i\v k\cdot\v x}\,f(\v k)\,,
\end{align}
where $f$ stands for any function or operator other than $a$ and $a^\dagger$, in particular for the colour charge density and for the coefficient functions introduced below. With these conventions
\begin{align}\label{eq:aadaggerX}
\big[a^a_i(k^+,\v z),a^{\dagger b}_j(p^+,\v z')\big]=2\pi\,\delta^{ab}\delta_{ij}\,\delta(k^+-p^+)\,\delta^{(2)}(\v z-\v z')\,.
\end{align}
We shall frequently need the Weizs\"acker--Williams field of a point charge and its Fourier transform,
\begin{align}\label{eq:kernels}
\Kv^i(\v x)\equiv\frac{x^i}{\v x^2}\,,\qquad \int_x e^{i\v k\cdot\v x}\,\Kv^i(\v x)=2\pi i\,\frac{k^i}{\v k^2}\,,\qquad
K(\v x,\v y,\v z)\equiv\Kv(\v x-\v z)\cdot\Kv(\v y-\v z)\,.
\end{align}

\subsection{Light-cone QCD Hamiltonian}\label{sec:hamiltonian}

The derivation of the QCD Hamiltonian in the light-cone gauge is standard~\cite{Brodsky:1997de,Lublinsky:2016meo}. We consider pure Yang--Mills theory, $N_f=0$. This is sufficient for the soft gluon sector at high energy; quarks can be added along the lines of~\cite{Lublinsky:2016meo}. After the elimination of $A^-$ the Hamiltonian is a functional of the transverse fields only and is the sum of an electric and a magnetic part,
\begin{align}\label{eq:HLC}
H=\int dx^-\,d^2x\,\Big[\frac12\,\Pi^a\,\Pi^a+\frac14\,F^a_{ij}\,F^a_{ij}\Big]\,,
\end{align}
with
\begin{align}\label{eq:PiF}
\Pi^a\equiv-\frac{1}{\partial^+}\big(D^{ab}_i\,\partial^+A^b_i\big)\,,\qquad
F^a_{ij}\equiv\partial_iA^a_j-\partial_jA^a_i-gf^{abc}A^b_iA^c_j\,,\qquad
D^{ab}_i\equiv\delta^{ab}\partial_i-gf^{abc}A^c_i\,.
\end{align}
Here $\partial^+\equiv\partial/\partial x^-$, $\partial_i\equiv\partial/\partial x^i$ and $g$ is the coupling. The expansion of Eq.~\eqref{eq:HLC} in powers of $g$, $H=H_0+H_{\rm int}$, and the quantisation of the transverse field are given in Appendix~\ref{app:hamiltonian}. The eigenstates of the free Hamiltonian $H_0$ are the Fock states of gluons with the energies~\eqref{eq:energy}; $H_{\rm int}$ contains the three- and four-gluon vertices and the instantaneous interactions characteristic of light-cone quantisation, Eq.~\eqref{eq:Hint}.

\subsection{Born--Oppenheimer formalism and the eikonal approximation}\label{sec:BO}

\paragraph{Valence and soft modes.} In the CGC description of a fast projectile~\cite{McLerran:1993ni,McLerran:1993ka,Jalilian-Marian:1997jx,Kovner:2005en} the gluon modes are divided according to their longitudinal momentum. The \emph{valence} modes have $k^+>\vee$, the \emph{soft} modes have
\begin{align}\label{eq:softwindow}
\Lambda<k^+<\vee\,,\qquad \vee\equiv\Lambda\,e^{\delta Y}\,.
\end{align}
Thus $\vee$ is the upper and $\Lambda$ the lower cut-off on the longitudinal momentum of the soft modes, and $\delta Y=\log(\vee/\Lambda)$ is the rapidity interval that they occupy; the notation is that of~\cite{Lublinsky:2016meo}. Modes with $k^+<\Lambda$ are not part of the projectile wave function considered here. The position of the boundary $\vee$ between valence and soft modes is arbitrary: a change of $\vee$ moves modes from one class to the other, and observables do not depend on it once the valence state is evolved accordingly. Accordingly the field is decomposed as
\begin{align}\label{eq:Asplit}
A^a_i(x)=\underline A^a_i(x)+\overline A^a_i(x)\,,
\end{align}
where the soft field $\underline A$ is given by Eq.~\eqref{eq:modeexp} with the $k^+$ integral restricted to $\Lambda<k^+<\vee$, and the valence field $\overline A$ by the same expression with $k^+>\vee$.

\paragraph{Born--Oppenheimer approximation.} The valence modes have much larger longitudinal momenta than the soft ones and, by Eq.~\eqref{eq:energy}, much smaller light-cone energies at a given transverse momentum. On the time scale of the soft modes they are frozen. The wave function of the soft modes is therefore found at a fixed configuration of the valence modes, which act as a static source~\cite{Kovner:2005en,Lublinsky:2016meo}. The only property of the valence modes that the soft modes are sensitive to, in the approximation described below, is their colour charge density integrated over $x^-$,
\begin{align}\label{eq:rhodef}
\rho^a(-\v p)=-if^{abc}\int_\vee^\infty\frac{dk^+}{2\pi}\int\frac{d^2k}{(2\pi)^2}\,a^{\dagger b}_j(k^+,\v k)\,a^c_j(k^+,\v k+\v p)\,.
\end{align}
It is an operator on the Hilbert space of the valence modes and commutes with the soft operators. It is Hermitian in coordinate space, so that $[\rho^a(\v p)]^\dagger=\rho^a(-\v p)$. From Eqs.~\eqref{eq:rhodef} and~\eqref{eq:aadagger} it follows that the charges at different colours do not commute but form a local $SU(N_c)$ algebra,
\begin{align}\label{eq:rhoalgebra}
\big[\rho^a(\v k),\rho^b(\v p)\big]=if^{abc}\,\rho^c(\v k+\v p)\,,\qquad
\big[\rho^a(\v x),\rho^b(\v y)\big]=if^{abc}\,\rho^c(\v x)\,\delta^{(2)}(\v x-\v y)\,.
\end{align}
We keep the charges as non-commuting operators throughout the paper. Products of charges therefore appear in a definite order, which has to be kept, and the commutator of two charges is one charge: a term with $n$ charges can be traded for a term with $n-1$ charges at the same order in the coupling.

\paragraph{Eikonal approximation.} When the decomposition~\eqref{eq:Asplit} is inserted into $H_{\rm int}$, terms with soft fields only, with valence fields only, and mixed terms are generated. We keep the leading power in the ratio of soft to valence longitudinal momenta, $k^+_{\rm soft}/k^+_{\rm valence}$. In practice this means that every term in which an inverse derivative $1/\partial^+$ acts on a product containing a valence field together with a soft field is dropped, and that the valence fields enter only through the bilinear $\overline A^b_j\,\partial^+\overline A^c_j$, whose integral over $x^-$ is the charge density~\eqref{eq:rhodef}. The terms that contain valence fields only do not act on the soft modes and are omitted. The interaction then is
\begin{align}\label{eq:Hdecomp}
H_{\rm int}\to H_g+H_{ggg}+H_{gg\text{-inst}}+H_{gggg}\,,
\end{align}
as in~\cite{Lublinsky:2016meo}; the four terms are given in Eqs.~\eqref{eq:Hpieces}--\eqref{eq:Hinst} of Appendix~\ref{app:hamiltonian}. $H_{ggg}$ is the three-gluon vertex of the soft modes, $H_{gg\text{-inst}}$ the instantaneous interaction of two soft gluons with the valence charge, and $H_{gggg}$ the four-gluon interaction of the soft modes, which is not needed in this paper. The term that drives the soft gluon wave function is the coupling of one soft gluon to the valence charge,
\begin{align}\label{eq:Hg}
H_g=\int_\Lambda^\vee\frac{dk^+}{2\pi}\int\frac{d^2k}{(2\pi)^2}\,\frac{g\,k^i}{\sqrt2\,(k^+)^{3/2}}\Big[a^{\dagger a}_i(k^+,\v k)\,\rho^a(-\v k)+a^a_i(k^+,\v k)\,\rho^a(\v k)\Big]\,.
\end{align}

\paragraph{Range of validity.} The expressions above are the leading terms of an expansion in $k^+_{\rm soft}/k^+_{\rm valence}$; power corrections in this ratio (non-eikonal corrections) are neglected. The charge density is not assumed to be small in this subsection. In the calculation of the wave function below we expand in powers of $g$ at fixed $\rho$, i.e.\ we consider a dilute projectile, $g\rho\ll1$; the order of a term in $g$ and the number of charges that it contains are then two separate labels.

\subsection{Light-cone quantization and Fock space expansion}\label{sec:fock}

\paragraph{States.} The soft vacuum $\ket0$ is annihilated by all soft annihilation operators. It is a state of the soft modes only: the state of the projectile is a vector in $\mathcal H_{\rm soft}\otimes\mathcal H_{\rm val}$, and $\ket0$ stands for $\ket0\otimes\ket v$ with $\ket v$ the unspecified state of the valence modes. The states with one and two soft gluons are $\ket{g^a_i(k)}=(2\pi)^{-3/2}a^{\dagger a}_i(k)\ket0$ and $\ket{g^a_i(k)\,g^b_j(p)}=(2\pi)^{-3}a^{\dagger a}_i(k)a^{\dagger b}_j(p)\ket0$, normalised to $\delta$-functions without factors of $2\pi$; in transverse coordinate space $\ket{g^a_i(k^+,\v z)}=(2\pi)^{-1/2}a^{\dagger a}_i(k^+,\v z)\ket0$. The complete definitions are collected in Appendix~\ref{app:states}.

\paragraph{Fock space expansion.} A state of the projectile is expanded in the Fock basis of the soft modes,
\begin{align}\label{eq:fock}
\ket\Psi=\psi_0\ket0+\int d^3k\ \psi^a_i(k)\ket{g^a_i(k)}+\frac12\int d^3k\,d^3p\ \psi^{ab}_{ij}(k,p)\ket{g^a_i(k)\,g^b_j(p)}+\dots
\end{align}
The expansion is in the soft modes only. The valence state is not expanded, and the coefficients $\psi_0,\psi^a_i,\dots$ are operators on $\mathcal H_{\rm val}$: they are functionals of the charge density $\rho$ that act on $\ket v$. In a product of such coefficients the order of the charges is kept, as required by Eq.~\eqref{eq:rhoalgebra}. A component of the wave function is characterised by two independent labels, the number $n$ of soft gluons and the number $m$ of charges that multiply it. At a given order in $g$ a component with a given $n$ receives contributions with different $m$; for instance the two-gluon component at order $g^2$ contains a term with one charge (a gluon emitted by the valence charge splits) and a term with two charges (two gluons are emitted). We label the components accordingly, e.g.\ $\ket{\Psi_{gg\rho}}$ and $\ket{\Psi_{gg\rho\rho}}$. Since by Eq.~\eqref{eq:rhoalgebra} a commutator of two charges is one charge, the label $m$ is defined only after a convention for the products of charges has been chosen. We use the decomposition
\begin{align}\label{eq:decomp}
\rho^a(\v k)\,\rho^b(\v p)=\frac12\big\{\rho^a(\v k),\rho^b(\v p)\big\}+\frac{i}{2}\,f^{abc}\,\rho^c(\v k+\v p)\,,
\end{align}
and count the anticommutator as a term with two charges and the commutator as a term with one charge. The anticommutator describes two emitters at different transverse positions; the commutator is local in coordinate space and is equivalent to one emitter. We refer to Eq.~\eqref{eq:decomp} as the decomposition procedure.

\subsection{Light-cone perturbation theory}\label{sec:lcpt}

Light-cone perturbation theory (LCPT) is time-independent (old-fashioned) perturbation theory for the Hamiltonian $H=H_0+V$, $V\equiv H_{\rm int}$, in the basis of the eigenstates of $H_0$~\cite{Lepage:1980fj,Brodsky:1997de}. The interaction acts on a free Fock state and generates the other Fock components; each insertion of $V$ is accompanied by an energy denominator, the difference between the energy of the initial state and the energy of the intermediate state that follows the insertion. In the present problem there is one difference to the textbook case. The matrix elements of $V$ between soft Fock states,
\begin{align}
V_{nm}\equiv\braket{n|H_{\rm int}|m}\,,
\end{align}
are not numbers but operators on the valence Hilbert space, since $H_g$ and $H_{gg\text{-inst}}$ contain $\rho$. They do not commute with each other. The energies $E_n$, on the other hand, are numbers, because the valence modes do not contribute to the energy in the approximation of Section~\ref{sec:BO}. The perturbative series therefore has the same form as for a number-valued perturbation, provided that the order of the matrix elements is kept.

\paragraph{Perturbation theory with non-commuting matrix elements.} Let $\ket i$ be an eigenstate of $H_0$ with energy $E_i$, and assume $V_{ii}=0$, which holds for all cases considered below. The state that develops from $\ket i$ when the interaction is switched on is, to second order,
\begin{align}\label{eq:pt2}
\ket{\Psi_i}=\ &\ket i\,\Big[1-\frac12\sum_{m\neq i}\frac{V_{im}V_{mi}}{(E_i-E_m)^2}\Big]
+\sum_{n\neq i}\ket n\,\frac{V_{ni}}{E_i-E_n}\notag\\
&+\sum_{n,m\neq i}\ket n\,\frac{V_{nm}\,V_{mi}}{(E_i-E_n)(E_i-E_m)}+O(V^3)\,.
\end{align}
The sums run over the free Fock states of the soft modes, and the later interaction stands to the left. The second and third terms are the Rayleigh--Schr\"odinger series; the term in the square bracket normalises the state. The corresponding series to third order for the soft vacuum, Eq.~\eqref{eq:pt3} of Appendix~\ref{app:pt}, which underlies the results quoted in Section~\ref{sec:vacin}, is the one of Appendix~E of~\cite{Lublinsky:2016meo}. The states are defined uniquely by two properties: $H$ does not connect dressed states of different energy, and the dressed states are orthonormal with a Hermitian component along the initial state (canonical normalisation). The second property fixes the order of the factors in the normalisation terms. We have verified the series, including the order of the factors, against an exact numerical construction (Appendix~\ref{app:pt}).

\paragraph{Energy denominators.} For the dressing of the soft vacuum the denominators are $E_0-E_n=-E_n<0$ and never vanish. For the dressing of a one-gluon state $\ket{g^a_i(k)}$ the denominators are $E_g(k)-E_n$. If the state $\ket n$ contains one gluon of momentum $p$, the denominator $E_g(k)-E_g(p)$ vanishes inside the integration region, on the surface $p^+\v k^2=k^+\v p^2$; if $\ket n$ contains two gluons that originate from a splitting of the gluon $k$, it is negative definite. We define the vanishing denominators as principal values,
\begin{align}\label{eq:PV}
\frac{1}{E_g(k)-E_g(p)}\to\PV\,\frac{1}{E_g(k)-E_g(p)}\,.
\end{align}
Eq.~\eqref{eq:PV} is part of the definition of the dressed states used in this paper, and all results below refer to it. It is the Hermitian part of the prescription $1/(E_i-E_n+i\varepsilon)$. The adiabatic prescription $E_i-E_n\to E_i-E_n+i\varepsilon$, which is used in~\cite{Radhakrishnan:2026soft}, differs from the principal value by terms proportional to $\delta(E_i-E_n)$, which have support on the energy shell; these terms are not included here. Vanishing denominators occur in one coefficient only, $\mf B_3$ of Section~\ref{sec:oneout}, and therefore in one row of the cross section (Section~\ref{sec:nuance_pv}).

\paragraph{Kinematic restrictions.} Every soft gluon, in the final state or in an intermediate state, has $\Lambda<p^+<\vee$. In addition, the longitudinal momentum that the last term of $H_{gg\text{-inst}}$ transfers between a soft gluon and the valence modes is bounded from below by $\Lambda$, $|p^+-q^+|>\Lambda$ in Eq.~\eqref{eq:Hinst}; the same restriction applies to the longitudinal momentum of a gluon that is emitted or absorbed by the valence charge through $H_g$, since this gluon is itself a soft mode. These restrictions fix the limits of all longitudinal integrals.

\paragraph{Matrix elements.} The matrix elements of $H_g$, $H_{ggg}$ and $H_{gg\text{-inst}}$ between states with up to two gluons, Eqs.~\eqref{eq:m1}--\eqref{eq:m5}, are listed in Appendix~\ref{app:matrixelements}. They coincide with those of~\cite{Lublinsky:2016meo}. All of them except the three-gluon vertex are operators on $\mathcal H_{\rm val}$.

\subsection{Soft gluon wave function in the CGC and the evolution operator \texorpdfstring{$\Omega$}{Omega}}\label{sec:omega_def}

\paragraph{Definition.} The states~\eqref{eq:pt2} are the dressed states of the soft modes in the presence of the valence charge. They are obtained from the free Fock states by an operator $\Omega$,
\begin{align}\label{eq:Omegadef}
\ket{\Psi_n}=\Omega\ket n\,,\qquad \Omega^\dagger\Omega=\Omega\,\Omega^\dagger=1\,,
\end{align}
which acts on $\mathcal H_{\rm soft}\otimes\mathcal H_{\rm val}$ and is unitary because it maps one orthonormal basis onto another. By the first of the two properties stated in Section~\ref{sec:lcpt}, $\Omega^\dagger H\Omega$ has no matrix elements between free soft Fock states of different energy: $\Omega$ diagonalises the Hamiltonian of the soft modes order by order in $g$~\cite{Kovner:2005en,Radhakrishnan:2026soft}. In particular
\begin{align}
\ket{\Psi_0}=\Omega\ket0
\end{align}
is the soft gluon wave function of the projectile: the cloud of soft gluons that accompanies the valence charges. When the rapidity interval of the soft modes is increased, $\Omega$ generates the gluons that populate the new interval, which is why it is called the evolution operator.

At the lowest order in $g$ the only interaction is $H_g$, and $\Omega=1+\int\frac{d^3p}{(2\pi)^3}\big[\mf A^b_j(p)\,a^{\dagger b}_j(p)-\mf A^{\dagger b}_j(p)\,a^b_j(p)\big]+O(g^2)$ with $\mf A^b_j(p)\propto g\,p^j\rho^b(-\v p)/\v p^2$ the Weizs\"acker--Williams field of the valence charge, Eq.~\eqref{eq:A1mom} below. For commuting charges this is the expansion of a coherent operator, which dresses the valence state with independently emitted gluons~\cite{Kovner:2005en}. Beyond the lowest order $\Omega$ contains the non-Abelian vertices and the commutators of the charges, and it is not a coherent operator.

\paragraph{Normal-ordered expansion.} Since any operator on the soft Fock space can be written as a sum of normal-ordered products of $a^\dagger$ and $a$, we parametrise
\begin{align}\label{eq:OmegaExp}
\Omega = \ &\mf{N}
+\int\frac{d^3p}{(2\pi)^3}\Big[\mf{A}^b_{1j}(p)\,a^b_j(p)+\mf{A}^b_{2j}(p)\,a^{\dagger b}_j(p)\Big]
\notag\\&
+\int\frac{d^3p}{(2\pi)^3}\frac{d^3q}{(2\pi)^3}\Big[\mf{B}^{cb}_{1kj}(p,q)\,a^b_j(p)a^c_k(q)+\mf{B}^{cb}_{2kj}(p,q)\,a^{\dagger b}_j(p)a^{\dagger c}_k(q)
\notag\\&\hspace{3.4cm}+\mf{B}^{cb}_{3kj}(p,q)\,a^{\dagger c}_k(q)a^b_j(p)\Big]
\notag\\&
+\int\frac{d^3p}{(2\pi)^3}\frac{d^3q}{(2\pi)^3}\Big[\mf{C}^{dcb}_{1lkj}(p,q)\,a^{\dagger d}_l(q)a^{\dagger c}_k(p-q)a^b_j(p)
\notag\\&\hspace{3.4cm}+\mf{C}^{dcb}_{2lkj}(p,q)\,a^{\dagger b}_j(p)a^c_k(p-q)a^d_l(q)\Big]+\dots
\end{align}
The coefficients $\mf N$, $\mf A_{1,2}$, $\mf B_{1,2,3}$, $\mf C_{1,2}$ are operators on $\mathcal H_{\rm val}$, i.e.\ functionals of $\rho$, and do not commute with each other; they commute with $a$ and $a^\dagger$. In each term the first index pair and the first momentum argument are tied to each other as indicated; note that in $\mf B^{cb}_{kj}(p,q)$ the \emph{second} colour and polarisation indices $(b,j)$ belong to the \emph{first} argument $p$. In the terms with three operators we have used that all vertices conserve the longitudinal and the transverse momentum of the soft gluons when $\rho$ is not involved; this is sufficient here because $\mf C_{1,2}$ are needed only at order $g$, where they are independent of $\rho$. The dots stand for terms with three operators that depend on $\rho$ and for terms with four or more operators. They are of order $g^2$ or higher and do not contribute to the single inclusive gluon cross section at the order considered in this paper (Section~\ref{sec:threeout}).

The products of operators in Eq.~\eqref{eq:OmegaExp} are symmetric under the exchange of two creation or of two annihilation operators, so that only the symmetric parts of the coefficients are defined. We impose
\begin{align}\label{eq:symm}
\mf B^{cb}_{1kj}(p,q)=\mf B^{bc}_{1jk}(q,p)\,,\qquad \mf B^{cb}_{2kj}(p,q)=\mf B^{bc}_{2jk}(q,p)\,,\qquad
\mf C^{dcb}_{1lkj}(p,q)=\mf C^{cdb}_{1klj}(p,p-q)\,,
\end{align}
and the same relation for $\mf C_2$.

\paragraph{Orders in the coupling.} In the expansion in $g$ at fixed $\rho$ the coefficients start at the following orders,
\begin{align}\label{eq:orders}
\mf N=1+\mf N_2\,,\quad \mf N_2=O(g^2)\,;\qquad
\mf A_{1,2}=\mf A^{(1)}_{1,2}+\mf A^{(3)}_{1,2}\,;\qquad
\mf B_{1,2,3}=O(g^2)\,;\qquad \mf C_{1,2}=O(g)\,,
\end{align}
where $\mf A^{(n)}=O(g^n)$. There are no terms of order $g^2$ in $\mf A_{1,2}$ and no terms of order $g$ in $\mf N$, because the interaction changes the number of gluons by one at order $g$ and the states are normalised canonically. When no confusion can arise we write $\mf A\equiv\mf A^{(1)}_2$ and $\mf C\equiv\mf C_1$.

\paragraph{Action on the vacuum and on a one-gluon state.} Acting with Eq.~\eqref{eq:OmegaExp} on the soft vacuum, only the terms without annihilation operators survive. With the normalisation~\eqref{eq:states},
\begin{align}\label{eq:Omegavac}
\Omega\ket0=\ &\mf N\ket0+(2\pi)^{3/2}\int\frac{d^3p}{(2\pi)^3}\,\mf A^b_{2j}(p)\ket{g^b_j(p)}\notag\\&+(2\pi)^3\int\frac{d^3p}{(2\pi)^3}\frac{d^3q}{(2\pi)^3}\,\mf B^{cb}_{2kj}(p,q)\ket{g^b_j(p)\,g^c_k(q)}+\dots
\end{align}
Acting on a one-gluon state, each annihilation operator has to be contracted with the gluon of the state, by Eq.~\eqref{eq:aadagger}, and one obtains
\begin{align}\label{eq:Omegag}
\Omega\ket{g^a_i(k)}=\ &\frac{1}{(2\pi)^{3/2}}\,\mf A^a_{1i}(k)\ket0+\mf N\ket{g^a_i(k)}+\int\frac{d^3p}{(2\pi)^3}\,\mf B^{ba}_{3ji}(k,p)\ket{g^b_j(p)}
\notag\\&
+(2\pi)^{3/2}\int\frac{d^3p}{(2\pi)^3}\,\mf A^b_{2j}(p)\ket{g^a_i(k)\,g^b_j(p)}
\notag\\&+(2\pi)^{3/2}\int\frac{d^3p}{(2\pi)^3}\,\mf C^{bca}_{1jki}(k,p)\ket{g^b_j(p)\,g^c_k(k-p)}+\dots,
\end{align}
where the dots are states with three or more gluons and terms of higher order. The terms have a simple interpretation. $\mf A_1$ absorbs the gluon; $\mf N$ and $\mf B_3$ leave the number of gluons unchanged, $\mf B_3$ describing the scattering of the gluon on the valence charges; $\mf A_2$ adds a gluon emitted by the valence charges to the incoming one, which is a spectator; $\mf C_1$ splits the gluon into two.

\paragraph{How the coefficients are determined.} The coefficients are fixed by two requirements:
\begin{itemize}
\item \emph{Matching.} $\Omega\ket n$ equals the state $\ket{\Psi_n}$ computed in LCPT, Eqs.~\eqref{eq:pt2} and~\eqref{eq:pt3}. Comparing Eq.~\eqref{eq:Omegavac} with $\ket{\Psi_0}$, Fock sector by Fock sector, fixes $\mf N$, $\mf A_2$ and $\mf B_2$. Comparing Eq.~\eqref{eq:Omegag} with the dressed one-gluon state fixes $\mf A_1$, $\mf B_3$ and $\mf C_1$, and reproduces $\mf N$ and $\mf A_2$.
\item \emph{Unitarity.} $\Omega^\dagger\Omega=1$, imposed order by order in $g$, relates the coefficients to each other. It determines the coefficients that multiply annihilation operators only, $\mf A_1$ and $\mf B_1$, and $\mf C_2$, in terms of those fixed by the matching, and it gives relations among the latter which serve as checks.
\end{itemize}
Each coefficient is probed by the transition between the Fock sectors that correspond to its annihilation and creation operators. It is therefore sufficient to consider the vacuum and the one-gluon state as initial states: all coefficients of Eq.~\eqref{eq:OmegaExp} have at most one annihilation operator, except $\mf B_1$ and $\mf C_2$, which follow from unitarity. The vacuum initial state was treated in~\cite{Lublinsky:2016meo}, and we quote the results in Section~\ref{sec:vacin}. The one-gluon initial state is treated in Section~\ref{sec:onegin}, and the unitarity relations are derived in Section~\ref{sec:unitarity}.
